\clearpage
\fancyhead[L]{\small COMET.}
\hypertarget{native-472-lower}{}
\begin{center}
\rule{0.55\linewidth}{0.4pt}\\[2pt]
Astronomische Nachrichten. Vol. XXII. No. 514. pp. 165. 166. 18 Aug. 1844.\\[2pt]
\rule{0.55\linewidth}{0.4pt}
\end{center}
\begin{center}\textsc{Gauss} to \textsc{Schumacher}. Göttingen, 18 August 1844.\end{center}

I am sending you a couple more comet observations. The weather has been consistently unfavourable, and even these scant observations were hard to obtain:

\begin{center}
\setlength{\tabcolsep}{4pt}
\begin{tabular}{l|c|c|c}
1844 & \shortstack{Mean time\\in Göttingen} & Right ascension & \shortstack{Northern\\declination} \\\hline
Aug. 11 & $9^{\rm u}\ 28^{\rm m}\ 16^{\rm s}$ & $211^\circ\ 34'\ 13''6$ & ---\quad ---\quad --- \\
 & $9\ 31\ 48$ & ---\quad ---\quad --- & $21^\circ\ 48'\ 12''6$ \\
17 & $8\ 34\ 23$ & $208\ 42\ 2.3$ & ---\quad ---\quad --- \\
 & $8\ 36\ 8$ & ---\quad ---\quad --- & $17\ 28\ 22.6$ \\
\end{tabular}
\end{center}

\Needspace{6\baselineskip}
\SourceLeaf{473}
On the 11th, I used three stars following the comet for comparison. The second will be the same star as the one listed in \textsc{Bessel}’s zone 460 as follows
\[
\begin{array}{c|c|c|c}
8.\,9 & 4 & 14^{\rm u}\ 7^{\rm m}\ 50^{\rm s}54 & 21^\circ\ 50'\ 21''9
\end{array}
\]
\noindent assuming that the observation was made at the third wire instead of the fourth and that the clock time was therefore printed with the wire interval mistakenly subtracted. My right ascension above rests on this hypothesis, which I have yet to confirm, and on the assumption that the reduction value \textsc{Bessel} used was $15^{\rm s}46$, or that the directly observed time of entry was $14^{\rm u}\ 8^{\rm m}\ 6^{\rm s}0$. If the second star were excluded altogether, the comet’s right ascension would be $4''2$ greater and would therefore have to be given as $211^\circ\ 34'\ 17''8$. Dr. \textsc{Goldschmidt}’s parabola still agrees quite well with the observations.

On 17 August, Piazzi XIII. 280 was used for comparison. The right ascension, however, comes from a single comparison.

